% Created 2026-08-20 Thu 15:45
% Intended LaTeX compiler: pdflatex
\documentclass[11pt]{article}
\usepackage[utf8]{inputenc}
\usepackage[T1]{fontenc}
\usepackage{graphicx}
\usepackage{longtable}
\usepackage{wrapfig}
\usepackage{rotating}
\usepackage[normalem]{ulem}
\usepackage{amsmath}
\usepackage{amssymb}
\usepackage{capt-of}
\usepackage{hyperref}
\usepackage[a4paper,margon=0.75in,top=0.85in]{geometry}
\usepackage{microtype}
\usepackage{parskip}
\usepackage{newpxtext}
\usepackage{newpxmath}
\usepackage{titlesec}
\titleformat{\section}{\large\bfseries}{}{0em}{}[{\titlerule[0.5pt]}]
\titleformat{\subsection}{\large\bfseries}{}{0em}{}
\usepackage{enumitem}
\usepackage{titling}
\usepackage{setspace}
\setstretch{1.15}
\setlist{noitemsep,topset=4pt}
\pretitle{\begin{center}\bfseries}
\posttitle{\end{center}\hr\vspace{1em}}
\usepackage{float}
\floatplacement{figure}{htbp}
\floatplacement{table}{htpb}
\setlength{\floatsep}{18pt plus 4pt minus 4pt}
\setlength{\textfloatsep}{20pt plus 4pt minus 4pt}
\usepackage{placeins}
\usepackage{booktabs}
\usepackage{graphicx}
\usepackage{caption}
\captionsetup{font=small, labelfont=bf}
\usepackage{hyperref}
\hypersetup{colorlinks=true, linkcolor=black, urlcolor=blue}
\author{Tyrone}
\date{}
\title{2026 Christmas Wishlist}
\hypersetup{
 pdfauthor={Tyrone},
 pdftitle={2026 Christmas Wishlist},
 pdfkeywords={},
 pdfsubject={},
 pdfcreator={Emacs 30.2 (Org mode 9.7.11)}, 
 pdflang={English}}
\begin{document}

\maketitle
\tableofcontents

Don't feel pressured to spend much. I'm happy with literally any of these.

\newpage
\section{Books}
\label{sec:org62a42d7}

\emph{The Priory of the Orange Tree} by Samantha Shannon

\emph{The Odyssey} Translation by McCrorie (2004) or Wilson (2017)

\emph{The Count of Monte Cristo}  by Alexandre Dumas and Auguste Maquet. Unabridged translation by Robin Buss (should be around 1,200+ pages).

\emph{East of Eden} by John Steinback

\emph{Small Things Like These} by Claire Keagan.

\emph{Goodnight, PunPun} Volume 2 by Inio Asano

\emph{One Hundred Years of Solitude} by Gabriel Garcia Marquez

\emph{The Fellowship of the Ring} by J. R. R. Tolkein

\newpage
\section{Stationery}
\label{sec:orgecb52db}

\subsection{Kokuyo Saxa Poche}
\label{sec:orge71a47a}
\begin{center}
\includegraphics[width=.9\linewidth]{/home/tyrone/Downloads/inks/saxa-poche.png}
\end{center}

I found some in \textbf{National Bookstore} for \textbf{399}.
\subsection{Pilot Kire-Na}
\label{sec:orgbc168aa}
\begin{center}
\includegraphics[width=.9\linewidth]{/home/tyrone/Downloads/inks/kire-na.jpg}
\end{center}

A set of 5 is just \textbf{\textasciitilde{}300}. I found some in both \textbf{Fully Booked} and \textbf{National Bookstore}.

\newpage
\section{Fountain Pen Inks}
\label{sec:org9a860b6}

Get me ink that comes in bottles or flasks, not those in small cartridges. I like the bottles. The cartridges may have proprietary designs, meaning they only fit in specific fountain pens from the same brand. But bottled ink can be used in any pen.

If the bottle advertises 'shading', that's good. I like ink that shifts colors.  Avoid inks that advertises 'shimmer', that's literally glitter. I normally don't like glittery inks. If it says 'sheen', that's okay. I don't really feel strongly about it.

For reference, \textasciitilde{}2 mL is enough ink for a month so 30mL would definitely last a full year.

All of these inks can be bought online from their official websites or in Lazada \& Shopee. You can find popular inks in large \textbf{Fully Booked} stores like in High Street or Mitsukoshi. Or in specialty stores like \textbf{Scribe} in SM Aura. You can also check \textbf{National Bookstore} but I'm pretty sure they don't have any of the premium inks.
\subsection{Specific Inks}
\label{sec:org1c8c043}

This is just a reference guide to give you an idea of what to expect, to check if you're getting scammed, or if you're getting a bad price. I'm just putting information about the  inks you'll likely find here so you don't confuse them and can tell the differences.

\textbf{You can pick any color, I don't mind}. Just pick what you think looks nice, you could get plain black or something more flamboyant. Just don't get me inks that are the same color like gifting me all blue inks or all purple inks.

The inks are listed in \textbf{descending order} of which I'd prefer.
\subsubsection{Pilot Iroshizuku}
\label{sec:orgbe159f8}
\begin{center}
\includegraphics[width=.9\linewidth]{/home/tyrone/Downloads/inks/iroshizuku.jpg}
\end{center}

Pilot's premium line of fountain pen ink. It has a core collection of 24 colors. I think they have the best looking bottles.They look like they're for perfume.

It comes in 50mL bottles, 15mL mini bottles, and 3mL cartridges. But the 15mL bottles don't really look very impressive. If you can, go for the 50mL bottle.

The ink itself is also very nice. I haven't seen any inks in an ugly color. They're also water-based so they write smoothly and don't clog fountain pens.

Estimated prices: \textasciitilde{}1,100 PHP for the 50mL bottles. Pretty good mL/PHP ratio. 50mL would last me two years if I use 2mL per month.

Suggested colors: \textbf{Tsuki-yo} or \textbf{Fuyu-syogun} or \textbf{Syo-ro} or \textbf{Sui-gyoku}. Like I said, they really don't have any bad colors.
\subsubsection{Sailor Shikiori}
\label{sec:orgc7cfd11}
\begin{center}
\includegraphics[width=.9\linewidth]{/home/tyrone/Downloads/inks/shikiori.png}
\end{center}

A premium, Japanese ink line themed around the four seasons. There's like 32 different colors. You can choose any color for me.

This line of ink generally comes in 10mL and 20mL glass bottles. 50mL used to be sold in round bottles but have been discontinued since 2019 I think. I'd prefer anything above 10mL. I don't think just 10mL is worth it.

Estimated price for a standard 20mL bottle: \textasciitilde{}600 PHP to \textasciitilde{}700 PHP. Those are still pretty poor volume-to-price ratios.

Suggested colors: \textbf{Yozakura} or \textbf{Tokiwa-matsu} or \textbf{Yamadori}.
\subsubsection{Sailor Ink Studio}
\label{sec:org92c0c29}
\begin{center}
\includegraphics[width=.9\linewidth]{/home/tyrone/Downloads/inks/ink_studio.png}
\end{center}

Technically, Sailor's most premium ink line. It is pretty niche and experimental. There's around 50 colors but the number changes a lot since colors are constantly getting added and discontinued.

Unlike other Sailor inks, each color is designated by three digits.

The first digit scales from 0 to 9 with 0 being extremely light, pastel shades to 9 being deep, almost black tones.

The second digit indicates color family: 2 is for Gray/Black, 3 is for red/pink, 4 is for the Blues, 5 is for purples, 6 is for greens, 7s and 8s are for  yellows, oranges, and browns.

The third digit is for specific undertones.

So \#123 is like a very light gray that has purple undertones.

They come only in 20mL glass bottles but those costs around \textasciitilde{}1,000 PHP. Not very cost-effective. I'd only recommend purchasing them if you they're on sale.

I think you can find standard colors for 675 in Scribe in SM Aura? But when I go they're usually out of stock.

Go with colors \textbf{\#123}, \textbf{\#224}, or \textbf{\#162}. In that order of preference. 
\subsubsection{Diamine}
\label{sec:orgee775cf}
\begin{center}
\includegraphics[width=.9\linewidth]{/home/tyrone/Downloads/inks/diamine.png}
\end{center}

A British brand with over a hundred colors. They're pretty affordable and the inks are reliable (they don't clog fountain pens easily).

A 30mL plastic bottle costs around 170 PHP to 220 PHP.

Large 80mL glass bottles costs around 550 PHP to 650 PHP. That's \textasciitilde{}8 pesos per mL or 16 pesos per month; good cost-effectiveness. 

But popular colors like Writer's Blood or Oxblood are usually out of stock.

Suggested colors: \textbf{Onyx Black} > \textbf{Writer's Blood} > \textbf{Oxford Blue} > \textbf{Ancient Copper}
\subsubsection{Lamy}
\label{sec:org8ce45b5}
\begin{center}
\includegraphics[width=.9\linewidth]{/home/tyrone/Downloads/inks/lamy.png}
\end{center}

Lamy ink comes in four main volumes: individual cartridges (labeled T10), plastic 30mL plastic bottles (T51), large 50mL glass bottles (T51), and high-end, crystal 30mL bottles (T53).

The standard colors come in Blue (which is erasable with water), Black, Blue-Black, Red, Green, Turquoise, and Violet.

T51 bottles will likely cost \textasciitilde{}500 PHP to \textasciitilde{}600 PHP.

T52 bottles are sold officially on Lazada for 669 PHP. Prices are also 669 on their official Lamy website.

T53 bottles are sold I think for \textasciitilde{}1,000 PHP.

The T52 bottles come in a few standard colors that are noted online to be muted and matte: \textbf{Turquoise} > \textbf{Petrol} > \textbf{Blue-Black} > \textbf{Dark Lilac} > \textbf{Black} > all the other colors.

While the T53 line have really saturated, intense dyed colors. The are themed after gemstones: \textbf{Amazonite} > \textbf{Agate} > \textbf{Obsidian} > \textbf{Benitoite} > all the others.
\subsubsection{Colorverse}
\label{sec:orgf870049}
\begin{center}
\includegraphics[width=.9\linewidth]{/home/tyrone/Downloads/inks/colorverse.png}
\end{center}

A premium South Korean brand famous for their inks themed around space and space. They are very difficult to find locally.

They come in 5mL mini bottles, 15mL teardrop-shaped bottles, 30mL oval bottles, 40mL bottles, 65 teardrop bottles, and combo sets of 65mL + 15mL bottles.

The 5mL bottles cost \textasciitilde{}500; that's crazy. That's 100 PHP per mL !

I found 30mL bottles at 1,200 PHP but selection is limited. I can't find any other sizes but I do see a lot of their empty bottles for sale. You \textbf{probably won't find any} but it'll be cool if you do.

Suggested colors: \textbf{Vervain} or \textbf{Terra}
\subsubsection{Sailor Manyo}
\label{sec:org53cdd22}
\begin{center}
\includegraphics[width=.9\linewidth]{/home/tyrone/Downloads/inks/manyo.jpeg}
\end{center}

A less premium line of ink from Sailor.It is themed around traditional Japanese poems. It has 35 main colors.

It comes in 50mL square glass bottles, 3mL ink samples, and 20mL bottles. But officially, the 20mL bottles comes in gift sets of 4 (totaling 80mL).

Estimated price range: +1,120 PHP for individual 50mL bottles (for standard colors) while +2,400 PHP for the 4-bottle sets (\textasciitilde{}600 per bottle).

I imagine you can find lone 20mL bottles from resellers.

Suggested colors: \textbf{Yamabuki} or \textbf{Kikyou}.
\subsubsection{Platinum}
\label{sec:orgb8515ce}
\begin{center}
\includegraphics[width=.9\linewidth]{/home/tyrone/Downloads/inks/platinum.png}
\end{center}

A top Japanese brands that offer different types of fountain pen inks.

There are \uline{Standard Dye} inks for everyday use. They come in cartridge packs that range from 130 PHP to 300 PHP while the 60mL bottles range from 750 PHP to 1,400 PHP depending on color and if they are limited editions.

There are \uline{Pigmented Inks}. They are for archival uses and water-resistant. That means they're more prone to clogging pens. They also come in cartridges or 60mL glass bottles. The glass bottles usually retail above 1,400 PHP.

There are the \uline{Classic Inks (Iron-Gall)}. They are interesting because the inks oxidizes and darkens as they dry. The 60mL ink retail at 1,400 to around 1,500.

There is the \uline{Mixable (Mix-Free)} Inks that are special because they can be mixed with each other safely, without causing unintended chemical reactions or clogging. They costs similarly to the higher-end Standard inks, so around 1,400 PHP per 60mL.

Lastly, there is the \uline{Chou Kuro Black}. This is the most premium ink line Platinum offers. It is called the "blackest of the black" inks. It's double the price of Standard Ink: 2,800 PHP.

If you choose to get me a Platinum ink, avoid the Mixable line because their whole point is to be mixed with another color but that would need at least two bottles (2 x 1,400 = 2,800). That's way too expensive.

I'd like to try \textbf{Pigmented Ink Carbon Black} if you can find any. \textbf{Classic Sepia Black} too would be interesting.
\subsubsection{Vinta}
\label{sec:org9e7164f}
\begin{center}
\includegraphics[width=.9\linewidth]{/home/tyrone/Downloads/inks/vinta.png}
\end{center}

A local, artisanal ink. Each color is themed around a major event in Filipino history. There's a ton of color but they're all shimmer heavy. I'm kind of interested. If you do get an ink, try aiming for subtler colors.

I think they only sell inks online, they don't sell in physical stores except through resellers.

Their 30mL bottles are like 449 to \textasciitilde{}600 PHP. 449 for the standard, shading inks while 599 for the shimmer-focused inks.
\subsubsection{Herbin}
\label{sec:orgcdd44de}
\begin{center}
\includegraphics[width=.9\linewidth]{/home/tyrone/Downloads/inks/herbin.jpg}
\end{center}

A really old French brand divided into two lines: \textbf{J. Herbin} and \textbf{Jacques Herbin}.

J. Herbin inks are meant to be affordable for everyday use while the Jacques Herbin inks are the premium line. The J. Herbin line comprises 35 core colors.

They specialize in shimmer inks.

They are sold in cartridges, 10mL \emph{Perle des Encres} jars, 30mL D-bottle Flacons, and 50mL cubes.

I found 30mL bottles for \textasciitilde{}700 PHP. I found a listing online selling 100mL bottles of standard ink for \textasciitilde{}1,300.

The 50mL bottles weirdly cost more with \textasciitilde{}1,500 PHP to \textasciitilde{}1,800 PHP.

Suggested colors: \textbf{Bleu Nuit} or \textbf{Vert Empire}; they lean less into the brand's signature shimmer.
\subsubsection{Asvine/Hongdian}
\label{sec:org9fab782}
\begin{center}
\includegraphics[width=.9\linewidth]{/home/tyrone/Downloads/inks/asvine-hongdian.jpg}
\end{center}

Both are Chinese brands owned by the same company. They are chemically similar and are very comparable with Jinhao inks, but are less messy, I guess.

Prices are slightly higher:

They are available in 60mL and 70mL bottles, depending on which brand. I found one in Lazada for around 400 PHP to under 500 PHP.

There's also 18mL bottles, I heard. But I can't find prices of them online.

As for colors, I guess \textbf{Peacock Blue} looks nice. It's up to you really.
\subsubsection{Jinhao}
\label{sec:org1f920bc}
\begin{center}
\includegraphics[width=.9\linewidth]{/home/tyrone/Downloads/inks/jinhao.jpg}
\end{center}

Jinhao is an entry-level, Chinese brand. The 30mL bottles should have a lot of colors to pick from. The larger 50mL bottles only have four main colors from what I know: Blue-Black > Black > Blue = Red. 

There are compact 30mL glass bottles sold in Lazada for \textasciitilde{}250 PHP.

There are also large 50mL glass bottles for under \textasciitilde{}300 PHP.

I heard this is a sleeper pick, surprisingly decent for it's price point.
\subsubsection{Pilot Standard}
\label{sec:org307afa8}
\begin{center}
\includegraphics[width=.9\linewidth]{/home/tyrone/Downloads/inks/pilot.png}
\end{center}

I already have these. They come in three main colors: Blue, Black, Blue-Black. I think there's red.

There's another line of Pilot inks called \textbf{Pilot Namiki}, they are the same chemically with Pilot Standard. So don't fall for listings saying Namiki is a premium ink.

Pilot Standard isn't labeled Pilot Standard. They're just called Pilot [color] like Pilot Blue or Pilot Black.

A 30mL bottle comes in at \textasciitilde{}200 PHP. It's what I normally buy. It is because of that, that I designate Pilot Standard as the baseline of fountain pen inks.

I prefer Blue-Black > Black > Blue >>> Red.
\end{document}
